\chapter{Le navigateur}

\section{Sources locales}
Le navigateur donne accès au stockage du téléphone et aux cartes SD. Sur Android~10 et plus, l'accès passe par le
sélecteur de dossiers du système : vous choisissez une fois le dossier de vos musiques, LHUNA s'en souvient.

\section{Catalogues en ligne}
LHUNA donne accès à plusieurs catalogues publics de musiques de puces, consultés en direct :
\begin{itemize}
  \item High Voltage SID Collection (Commodore~64) ;
  \item The Mod Archive, Modland et Amiga Music Preservation (modules de tracker, Amiga) ;
  \item ZXArt et ZXTunes (ZX Spectrum), Aygor ;
  \item Atari SAP Music Archive (Atari 8 bits) ;
  \item VGMRips (enregistrements VGM de consoles et bornes d'arcade) ;
  \item Joshw (rips de consoles), Khinsider, OverClocked ReMix, Aminet, Scene.org ;
  \item des radios en ligne.
\end{itemize}
\arediger{Description de chaque catalogue : contenu, taille, mainteneurs, conditions d'utilisation. À vérifier sur
leurs sites.}

\section{Archives et fichiers multi-morceaux}
LHUNA ouvre les archives (ZIP, RAR, 7z, LHA\dots) et les images de disquettes (TR-DOS, SCL\dots) comme des dossiers.
Un fichier qui contient plusieurs morceaux (SID, NSF, AY, KSS\dots) est en revanche lu directement : ses
sous-morceaux s'enchaînent, puis la lecture continue avec le fichier suivant du dossier.

\section{Recherche et sélection}
La loupe filtre le dossier affiché. Un appui long sur un élément démarre la sélection multiple ; la barre du haut
propose alors :
\begin{description}
  \item[Ajouter] demande la liste de destination et y ajoute les fichiers (les dossiers sont parcourus en entier) ;
  \item[Jouer ensuite] place les éléments sélectionnés dans la file d'attente (section~\ref{sec:file}).
\end{description}
